\section{Related work} 
\subsection{Data selection}
\paragraph{Traditional data selection methods and related concepts.}  
Data selection methods aim to identify subsets of training samples that maximize informativeness or diversity and minimize redundancy~\cite{ds_albalak2022survey}. Approaches leveraging submodular optimization and clustering have shown success in various tasks~\citep{ds_submodular_mirzasoleiman2016fast, ds_cluster_sener2018active, ds_submodular_wei2015submodularity}. Coreset selection techniques specifically construct small representative subsets that retain essential properties of the full dataset, often using importance weights or gradient similarity~\citep{ds_killamsetty2021grad,ds_cluster_sener2018active}. Learning-based methods integrate selection into training through training dynamics analysis~\citep{ds_paul2021deep,ds_Toneva2018AnES}, bilevel optimization~\citep{ds_bi_Killamsetty2020GLISTERGB,ds_bi_zhou2022probabilistic}, gradient properties~\citep{ds_sample_Johnson2018TrainingDM,ds_sample_Katharopoulos2018NotAS}, or uncertainty measures~\citep{ds_uncertain_Ash2019DeepBA}. Dataset distillation methods construct synthetic examples to compress large datasets, enabling faster or more efficient training by matching gradient statistics~\citep{ds_distill_Cazenavette2022DatasetDB,ds_distill_Deng2022RememberTP,ds_distill_wang2018dataset,ds_distill_Yu2023DatasetDA}. While highly compact, distilled sets often struggle to match real-data diversity or scale with heterogeneous corpora characteristic of large-scale model training. Curriculum learning~\citep{cl_bengio2009curriculum} dynamically orders or selects data according to difficulty, often progressing from easy to challenging samples~\citep{cl_hacohen2019power, cl_Soviany2021CurriculumLA}. A related body of work uses importance sampling and adaptive reweighting to prioritize data points that contribute most to model improvement or convergence~\citep{is_alain2015variance, cl_bengio2009curriculum,is_Loshchilov2015OnlineBS}, demonstrating efficiency gains especially in large-scale or imbalanced data regimes.

\paragraph{Data selection for LLM fine-tuning.}
Unlike traditional methods, data selection for LLMs often employs simpler heuristic rules due to the unprecedented scale of both models and datasets.
Efficient subset selection has become crucial in LLM fine-tuning~\cite{ds_albalak2022survey,Marone2023DataPR,Oren2019DistributionallyRL,Rae2021ScalingLM}. LIMA demonstrated strong performance with just 1,000 human-curated examples~\cite{ds_ift_Zhou2023LIMALI}, while automated approaches have emerged to replace manual curation. These include natural language indicators with BlendSearch~\cite{ds_ift_Cao2023InstructionMH}, semantic intention tagging~\citep{ds_ift_Lu2023InsTagIT}, model-based quality filtering~\citep{ds_ift_chen2023alpagasus,ds_ift_lian2023slimorca}, instruction difficulty metrics~\citep{ds_ift_Li2023FromQT}, and uncertainty-driven active learning~\citep{ds_ift_Kung2023ActiveIT}. Recent advances explore gradient-based techniques, including clustered coreset selection~\citep{ds_ift_cluster_Zhang2024TAGCOSTG} and graph methods~\citep{ds_ift_Zhao2025BeyondSA}, alongside simple yet effective heuristics like length-based prioritization~\citep{ds_ift_Zhao2024LongIM}. Novel approaches leverage small model training trajectories to guide selection for larger models~\citep{ds_ift_Yang2024SmallToLargeS} and sparse autoencoders for diversity-driven selection~\citep{ds_ift_Yang2025DiversitydrivenDS}. DEITA provides a unified framework balancing quality, complexity, and diversity~\citep{ds_ift_Liu2023WhatMG}. These methods achieve comparable or superior performance to full-dataset training while reducing computational requirements, demonstrating that strategically selected subsets can match or exceed naive scaling.

\paragraph{The compute-aware gap.}
Existing methods typically operate with predefined data budgets without considering computational constraints. Recent studies~\citep{ds_ift_Diddee2024ChasingRI,ds_ift_Shen2024RethinkingDS,ds_ift_Xia2024RethinkingDS} reveal that sophisticated selection strategies often fail to consistently outperform random selection across varied experimental settings. \cite{Yin2024ComputeConstrainedDS} further demonstrate that when computational budgets are explicitly factored in, previously effective data selection approaches rarely remain optimal. This computational perspective motivates our approach to data selection that explicitly accounts for training budget constraints.

\subsection{Bilevel optimization}
\paragraph{Algorithmic advances}
Early work on bilevel optimization in deep learning focused on implicit differentiation and Hessian–vector products~\citep{bi_domke2012generic,bi_grazzi2020iteration,bi_pedregosa2016hyperparameter} as well as iterative differentiation with truncated back‑propagation~\citep{bi_maclaurin2015gradient,bi_shaban2019truncated, bi_zhou2022model}. These methods scale poorly because they require nested loops and second‑order information.  Stochastic variants~\citep{bi_Chen2021ASM,bi_Ghadimi2018ApproximationMF,bi_Hong2020ATF,bi_Ji2020BilevelOC,bi_Khanduri2021ANA} cut per‑iteration cost but still depend on Jacobian/Hessian evaluations.  Finite‑difference estimators~\citep{bi_Sow2021OnTC,bi_Yang2023AchievingOC} remove explicit Hessians yet introduce instability.  Most recently, fully first‑order, penalty‑based formulations~\citep{bi_Chen2023NearOptimalNB,kwon2023fully,bi_Ye2022BOMEBO} recast the inner optimality conditions as constraints, eliminating higher‑order derivatives and enabling training with only standard gradients. This advancement enables bilevel algorithms to be applied to increasingly larger‑scale problems.   

\paragraph{Practical applications}  
Bilevel optimization has become a fundamental tool in machine learning, initially enabling hyperparameter optimization by jointly optimizing model parameters and hyperparameters~\citep{bi_app_Franceschi2018BilevelPF}. Key applications include learning-rate and weight-decay tuning~\citep{bi_franceschi2017forward,bi_lorraine2020optimizing}. In neural architecture search, DARTS frames architecture parameters as outer variables and weights as inner variables, achieving efficient architecture discovery~\citep{bi_app_Liu2018DARTSDA}. Resource-aware tasks benefit from probabilistic bilevel formulations for memory-constrained coreset selection~\citep{ds_bi_zhou2022probabilistic}, and gradient-matching methods for learning compact synthetic datasets~\citep{ds_distill_wang2018dataset}. Extending to large-scale settings, ScaleBiO applies first-order bilevel data reweighting at the scale of billions of tokens in language modeling, demonstrating improved performance with manageable computational overhead~\citep{pan2024scalebio}. These developments highlight bilevel optimization’s versatility in structuring complex learning problems, enabling efficient parameter tuning, architecture search, and data management, and suggest continued expansion into broader applications.   

